\documentclass[a4paper]{article}

\usepackage[fontset=fandol]{ctex}
\usepackage{amsmath, amssymb}
\usepackage{graphicx}
\usepackage[margin=2.5cm]{geometry}
\usepackage[most]{tcolorbox}
\usepackage{listings}
\usepackage{hyperref}
\usepackage{booktabs}
\usepackage{subcaption}
\usepackage{float}
\usepackage{tikz}

\setlength{\parindent}{2em}
\setlength{\parskip}{0.35em}
\sloppy

\newtcolorbox{knowledgebox}[1]{
    enhanced, colback=blue!5!white, colframe=blue!75!black, colbacktitle=blue!75!black,
    coltitle=white, fonttitle=\bfseries, title={#1}, attach boxed title to top left={yshift=-2mm, xshift=2mm},
    boxrule=1pt, sharp corners
}

\newtcolorbox{importantbox}[1]{
    enhanced, colback=yellow!10!white, colframe=yellow!80!black, colbacktitle=yellow!80!black,
    coltitle=black, fonttitle=\bfseries, title={#1}, sharp corners
}

\newtcolorbox{warningbox}[1]{
    enhanced, colback=red!5!white, colframe=red!75!black, colbacktitle=red!75!black,
    coltitle=white, fonttitle=\bfseries, title={#1}, sharp corners
}

\lstset{
    basicstyle=\ttfamily\small,
    keywordstyle=\color{blue},
    stringstyle=\color{red!60!black},
    commentstyle=\color{green!60!black},
    breaklines=true,
    frame=single,
    numbers=left,
    numberstyle=\tiny\color{gray},
    captionpos=b,
    extendedchars=false
}

\DeclareMathOperator*{\argmin}{arg\,min}
\newcommand{\HH}{\mathcal{H}}
\newcommand{\R}{\mathbb{R}}

\newcommand{\notetitle}{卷积神经网络 CNN}
\newcommand{\notesubtitle}{李宏毅《机器学习》2025 · Network Architecture Designed for Image}
\newcommand{\noteauthors}{基于李宏毅老师课程整理}
\newcommand{\notedate}{\today}
\newcommand{\videochannel}{李宏毅深度学习课堂（Bilibili）}
\newcommand{\videoduration}{约 56 分钟}
\newcommand{\videourl}{https://www.bilibili.com/video/BV1TAtwzTE1S?p=9}

\begin{document}
\pagestyle{empty}

\begin{center}
    \vspace*{1.5cm}
    {\Huge\bfseries \notetitle\par}
    \vspace{0.4cm}
    {\LARGE \notesubtitle\par}
    \vspace{0.8cm}
    \includegraphics[width=0.62\textwidth]{video.jpg}\par
    \vspace{0.8cm}
    {\large \noteauthors\par}
    {\large \notedate\par}
    \vspace{0.8cm}
    \begin{tcolorbox}[width=0.82\textwidth,center,sharp corners,
                      colback=black!3!white,colframe=black!50!white]
        \centering
        \textbf{视频信息}\\[3pt]
        频道：\videochannel \\
        时长：\videoduration \\
        链接：\href{\videourl}{\videourl}
    \end{tcolorbox}
\end{center}

\newpage
\tableofcontents
\newpage
\pagestyle{plain}

% =====================================================================
\section{从影像分类看架构设计}
% =====================================================================

这一讲开始讨论 network architecture 的设计。CNN 的重点不只是「有一种叫 convolution 的运算」，而是展示一个更一般的思想：\textbf{如果任务本身有结构，我们可以把这些结构写进网络架构里，让模型少学一些不必要的自由度。}

本讲的任务是 image classification：输入一张图片，输出图片中物件所属类别。若类别有 $K$ 种，目标标签通常写成 $K$ 维 one-hot vector；网络输出经过 softmax 后得到分布 $\boldsymbol{y}'$，训练时让 $\boldsymbol{y}'$ 与真实标签 $\hat{\boldsymbol{y}}$ 的 cross-entropy 越小越好。

\begin{figure}[H]
    \centering
    \includegraphics[width=0.86\textwidth]{figures/fig01_image_classification_task.jpg}
    \caption{影像分类任务：输入固定大小图片，输出类别。画面依据：约 01:45--03:10，字幕说明输入图像大小固定、输出用 one-hot label 与 softmax/cross-entropy。}
    \label{fig:image_task}
\end{figure}

在这堂课里，老师一直把 CNN 放在 P08 的泛化观点下理解：fully connected network 的函数集合很大，表达弹性很高，但参数多、候选函数多，也更容易在有限资料上 overfit。CNN 则把一些影像先验写死在架构里，牺牲一部分通用性，换来更适合影像任务的 inductive bias。

\begin{importantbox}{CNN 的核心不是更万能，而是更有偏置}
    Fully connected layer 可以学习任何像素间的连接关系；CNN 则预设「重要 pattern 通常是局部的」以及「同一个 pattern 可以出现在不同位置」。这些预设如果符合任务，就能减少参数、降低过拟合风险，并让模型更容易学到有用特征；如果任务不符合这些预设，CNN 反而可能不是好选择。
\end{importantbox}

% =====================================================================
\section{图片作为输入：三维 Tensor 与全连接的代价}
% =====================================================================

\subsection{图片在机器眼中是什么？}

一张彩色图片对人来说是视觉画面，对机器来说是数字阵列。若图片大小为 $100\times 100$，并且有 RGB 三个 channel，它就是一个
\[
    100\times 100\times 3
\]
的三维 tensor。前两个维度是空间位置，高和宽；第三个维度是颜色 channel。每个 pixel 在每个 channel 上都有一个强度数值。

\begin{figure}[H]
    \centering
    \includegraphics[width=0.86\textwidth]{figures/fig02_image_as_tensor.jpg}
    \caption{彩色图片可看成三维 tensor：宽、高与 RGB 三个 channel；若拉直，则变成长度 $100\times100\times3$ 的向量。画面依据：约 03:15--05:40。}
    \label{fig:image_tensor}
\end{figure}

到目前为止，课程中最常见的输入是向量。最直觉的做法就是把图片拉直：
\[
    x\in\R^{100\times100\times3}=\R^{30000}.
\]
这样就可以把图片丢进 fully connected network。这个做法在数学上没有问题，但在参数量上很快变得昂贵。

\subsection{全连接网络为什么太贵？}

假设第一层有 $1000$ 个 neuron。每个 neuron 都要对输入向量的每一维放一个 weight，因此第一层权重数约为
\[
    1000\times100\times100\times3
    =
    3\times10^7.
\]
这还只是第一层。若图片更大、第一层 neuron 更多，参数量会继续膨胀。

\begin{figure}[H]
    \centering
    \includegraphics[width=0.86\textwidth]{figures/fig03_fully_connected_cost.jpg}
    \caption{把 $100\times100\times3$ 图片拉直后接到 1000 个 neuron，第一层就有约 $3\times10^7$ 个 weight。画面依据：约 05:45--06:35。}
    \label{fig:fc_cost}
\end{figure}

参数多并不自动等于不好。参数多通常意味着模型更有弹性，理想 loss 可能更低；但在资料有限时，模型也更容易选择到只适合训练集的函数。CNN 要做的事，就是利用影像任务本身的结构，删除一些不必要的连接关系。

\begin{warningbox}{不要把参数减少理解成单纯省内存}
    CNN 当然能减少计算和参数，但这不是全部意义。更重要的是它改变了模型家族：它不再允许每个 neuron 任意看全图，也不再允许同一类 pattern 在不同位置各自学一套完全独立参数。这些约束让模型更偏向影像问题中常见的规则。
\end{warningbox}

% =====================================================================
\section{观察一：重要 Pattern 往往是局部的}
% =====================================================================

在判断一张图里有没有鸟时，人并不一定逐像素理解整张图。我们往往看见鸟嘴、眼睛、爪子、羽毛等关键 pattern，就做出判断。机器也可以有类似设计：某些 neuron 负责侦测小 pattern，而小 pattern 不需要看完整张图。

例如要判断某个小区域有没有鸟嘴，通常看局部范围就够了；不需要把整张 $100\times100$ 图都输入给这个 neuron。因此 CNN 的第一项简化是 \textbf{receptive field}。

\begin{figure}[H]
    \centering
    \includegraphics[width=0.86\textwidth]{figures/fig04_receptive_field_basic.jpg}
    \caption{Receptive field：一个 neuron 只看图片中的局部区域，而不是整张图。画面依据：约 10:45--11:55。}
    \label{fig:rf_basic}
\end{figure}

若 receptive field 大小为 $3\times3$，且覆盖全部 3 个 channel，则一个 neuron 的输入只有
\[
    3\times3\times3=27
\]
个数值。它只需要 27 个 weight 加一个 bias，而不是对整张图的 30000 个数值都放权重。

典型设定包括：
\begin{itemize}
    \item \textbf{all channels}：一般影像 CNN 的 receptive field 会覆盖所有输入 channel。
    \item \textbf{kernel size}：只描述空间大小，如 $3\times3$、$5\times5$。
    \item \textbf{stride}：receptive field 在图片上移动的步长，常见为 1 或 2。
    \item \textbf{padding}：当 receptive field 超出边界时补值，最常见是 zero padding。
    \item \textbf{multiple neurons per field}：同一个局部区域通常由一组 neuron 处理，如 64 个。
\end{itemize}

\begin{figure}[H]
    \centering
    \includegraphics[width=0.86\textwidth]{figures/fig05_receptive_field_typical.jpg}
    \caption{典型 receptive field 设置：覆盖所有 channel，kernel size 如 $3\times3$，stride 决定移动间隔，每个 receptive field 可有一组 neuron。画面依据：约 15:45--18:10。}
    \label{fig:rf_typical}
\end{figure}

\begin{knowledgebox}{为什么小 kernel 也能看见大 pattern？}
    单层 $3\times3$ filter 只看局部，但多层 convolution 堆叠后，后一层的一个位置对应到原图中更大的范围。例如第二层的 $3\times3$ 区域，其每个值已经来自第一层的局部卷积结果，因此映回原图时可能对应 $5\times5$ 或更大区域。网络越深，有效 receptive field 越大。
\end{knowledgebox}

% =====================================================================
\section{观察二：同样 Pattern 可以出现在不同位置}
% =====================================================================

第二个影像先验是：同一个 pattern 可以出现在图片的不同区域。鸟嘴可能在左上角，也可能在中间；但「侦测鸟嘴」这件事本身不应该因为位置不同就完全重学一套参数。

若每个位置都有独立 neuron 去学「鸟嘴侦测器」，参数会大量重复。CNN 因此引入 \textbf{parameter sharing}：不同位置的 neuron 可以共享同一组 weight，只是它们看的输入区域不同。

\begin{figure}[H]
    \centering
    \includegraphics[width=0.86\textwidth]{figures/fig06_parameter_sharing.jpg}
    \caption{Parameter sharing：不同 receptive field 中的 neuron 可使用同一组 weight；参数相同但输入区域不同，因此输出不必相同。画面依据：约 22:45--24:30。}
    \label{fig:param_sharing}
\end{figure}

用公式写，如果两个位置的局部输入分别是 $\boldsymbol{x}$ 与 $\boldsymbol{x}'$，共享参数为 $\boldsymbol{w}$，则输出可写成
\[
    a=\sigma(\boldsymbol{w}^{T}\boldsymbol{x}+b),
    \qquad
    a'=\sigma(\boldsymbol{w}^{T}\boldsymbol{x}'+b).
\]
虽然 $\boldsymbol{w}$ 和 $b$ 相同，只要 $\boldsymbol{x}\ne\boldsymbol{x}'$，输出通常也不同。共享参数不是让所有位置输出一样，而是让同一种侦测器在所有位置复用。

\begin{importantbox}{Convolutional layer = receptive field + parameter sharing}
    Receptive field 限制每个 neuron 只看局部；parameter sharing 限制不同位置的同类 neuron 共用参数。两者合起来，就是 convolutional layer 的核心。它比 fully connected layer 更受限制，model bias 更大，但这个 bias 正是为影像设计的。
\end{importantbox}

\begin{figure}[H]
    \centering
    \includegraphics[width=0.86\textwidth]{figures/fig07_convolutional_layer_bias.jpg}
    \caption{Convolutional layer 是 fully connected layer 的受限版本：先加入 receptive field，再加入 parameter sharing。画面依据：约 26:20--28:55。}
    \label{fig:conv_bias}
\end{figure}

这也解释了为什么 CNN 不是“比较强的全连接层”。它从表达能力上其实更窄，但窄在正确方向上：它把影像任务里不太需要的自由度拿掉。

% =====================================================================
\section{从 Filter 的故事看 Convolution}
% =====================================================================

同一件事也可以用 filter 来讲。一个 convolutional layer 里有一组 filters。每个 filter 是一个小 tensor，例如 $3\times3\times C$，其中 $C$ 是输入 channel 数。Filter 的参数要通过训练学出来；训练完成后，一个 filter 可以理解为某类局部 pattern 的侦测器。

\begin{figure}[H]
    \centering
    \includegraphics[width=0.86\textwidth]{figures/fig08_filters_story.jpg}
    \caption{Filter 版本的 CNN 叙事：每个 filter 是 $3\times3\times channel$ 的小 tensor，用来在图片中侦测某种局部 pattern。画面依据：约 29:40--30:35。}
    \label{fig:filters}
\end{figure}

卷积运算可以理解成：把 filter 放到图片某个位置，取局部 patch 与 filter 的 inner product，再把 filter 依照 stride 滑到下一个位置重复计算。

若输入 patch 为 $X_{i:i+2,j:j+2}$，filter 为 $W\in\R^{3\times3}$，则单个输出可写成
\[
    z_{i,j}
    =
    \sum_{u=0}^{2}\sum_{v=0}^{2}
    W_{u,v}X_{i+u,j+v}
    + b.
\]
彩色图片或多 channel feature map 时，还要再对 channel 维度求和。

\begin{figure}[H]
    \centering
    \includegraphics[width=0.86\textwidth]{figures/fig09_convolution_numeric.jpg}
    \caption{卷积的数值例子：filter 与局部 patch 做 inner product；filter 扫过整张图后得到一张输出图。画面依据：约 31:30--33:20。}
    \label{fig:conv_numeric}
\end{figure}

一个 filter 扫完整张图会得到一张二维输出；多个 filter 会得到多张输出。这些输出合在一起叫 \textbf{feature map}。若一层有 64 个 filters，输出就有 64 个 channel。下一层 convolution 的 filter 深度必须等于前一层输出 channel 数。

例如第一层输入是 RGB 图片，filter 形状可能是 $3\times3\times3$；若第一层有 64 个 filters，输出 feature map 有 64 个 channel；第二层 filter 就可能是 $3\times3\times64$。

\begin{figure}[H]
    \centering
    \includegraphics[width=0.86\textwidth]{figures/fig10_two_stories_compare.jpg}
    \caption{Neuron 版本与 filter 版本是同一件事：共享参数的 neuron，对应到一个 filter 扫过整张图。画面依据：约 37:40--39:25。}
    \label{fig:two_stories}
\end{figure}

\begin{knowledgebox}{为什么叫 convolution？}
    从 neuron 视角看，是多个不同位置的 neuron 共用同一组参数；从 filter 视角看，是同一个 filter 在整张图片上滑动计算。后一种操作就是 convolutional layer 名字里的 convolution。
\end{knowledgebox}

% =====================================================================
\section{Pooling：下采样与资讯取舍}
% =====================================================================

除了 convolution，经典 CNN 常搭配 pooling。Pooling 来自第三个观察：把图片做适度 subsampling，通常不会改变图片里的物件。例如一张鸟的图片缩小后还是鸟。

\begin{figure}[H]
    \centering
    \includegraphics[width=0.86\textwidth]{figures/fig11_subsampling_observation.jpg}
    \caption{Pooling 的直觉来源：适度缩小图片不改变物件类别。画面依据：约 40:10--40:35。}
    \label{fig:subsampling}
\end{figure}

Pooling 本身没有可学习参数，更像一个固定 operator。以 max pooling 为例，把 feature map 划成若干小块，如 $2\times2$，每块取最大值作为代表：
\[
    y_{i,j}
    =
    \max_{(u,v)\in\Omega_{i,j}} x_{u,v}.
\]
这样空间大小变小，channel 数不变。例如 $4\times4\times64$ 经过 $2\times2$ max pooling 后可能变成 $2\times2\times64$。

\begin{figure}[H]
    \centering
    \includegraphics[width=0.86\textwidth]{figures/fig12_pooling_result.jpg}
    \caption{Convolution 后得到多 channel feature map；pooling 会降低空间解析度，但通常保留 channel 数。画面依据：约 41:10--42:45。}
    \label{fig:pooling}
\end{figure}

\begin{warningbox}{Pooling 不是永远该用}
    Pooling 会减少计算量，也能带来一定程度的位置稳定性；但它也丢掉资讯。若任务非常依赖精细空间位置，例如围棋，随便下采样可能改变问题本身。近年的很多视觉架构也会减少或取消 pooling，改用更深的 convolution 或其他下采样设计。
\end{warningbox}

% =====================================================================
\section{完整 CNN：Convolution、Pooling、Flatten、Fully Connected}
% =====================================================================

经典 CNN 的整体结构通常是：
\[
    \text{Image}
    \rightarrow
    \text{Convolution}
    \rightarrow
    \text{Pooling}
    \rightarrow
    \cdots
    \rightarrow
    \text{Flatten}
    \rightarrow
    \text{Fully Connected}
    \rightarrow
    \text{Softmax}.
\]
Convolution/Pooling 负责把原始像素转成越来越抽象的 feature map；Flatten 把矩阵状 feature map 拉直成向量；fully connected layers 再根据这些高层特征做最终分类。

\begin{figure}[H]
    \centering
    \includegraphics[width=0.86\textwidth]{figures/fig13_whole_cnn.jpg}
    \caption{完整 CNN：多次 convolution/pooling 后 flatten，再接 fully connected layers 和 softmax 输出类别。画面依据：约 43:55--44:45。}
    \label{fig:whole_cnn}
\end{figure}

这里可以把前面的三种设计合起来看：
\begin{itemize}
    \item Locality：filter 只看局部，减少单个侦测器输入维度。
    \item Parameter sharing：同一个 filter 扫过全图，减少重复学习。
    \item Hierarchy：多层 convolution 让低层 pattern 组合成高层 pattern。
\end{itemize}

% =====================================================================
\section{应用例子：为什么 CNN 可以下围棋？}
% =====================================================================

CNN 不只用于自然图片，也曾用于 AlphaGo 的围棋模型。围棋可以写成分类任务：输入棋盘状态，输出下一步落子位置。棋盘有 $19\times19$ 个位置，因此下一步位置可看成 $19\times19$ 个类别之一。

棋盘也可以看成一张小图片：空间解析度是 $19\times19$，每个位置的 channel 描述该点状态，例如黑子、白子、空位、是否被叫吃等。在 AlphaGo 的论文描述中，输入是
\[
    19\times19\times48
\]
的 image stack。

\begin{figure}[H]
    \centering
    \includegraphics[width=0.86\textwidth]{figures/fig14_go_as_image.jpg}
    \caption{围棋可转成影像式输入：$19\times19$ 棋盘位置加上多组 feature planes，输出是 $19\times19$ 个可能落子位置之一。画面依据：约 45:05--48:15。}
    \label{fig:go_image}
\end{figure}

为什么围棋适合 CNN？因为它和影像有两个相似结构：
\begin{itemize}
    \item 局部 pattern 重要：某些棋形只需要看小范围就能判断，例如叫吃。
    \item 同样 pattern 可在不同位置出现：同一种棋形可能出现在棋盘任何角落。
\end{itemize}

\begin{figure}[H]
    \centering
    \includegraphics[width=0.86\textwidth]{figures/fig15_why_cnn_for_go.jpg}
    \caption{围棋与影像共享 CNN 的两项先验：重要 pattern 通常局部，同样 pattern 可出现在不同区域。画面依据：约 48:20--49:45。}
    \label{fig:go_why}
\end{figure}

但围棋不适合 pooling。自然图片缩小一点仍是同一物件，棋盘却不能随便删掉某些行列；一个位置消失，棋局就变了。老师特意查 AlphaGo 的论文附录，指出其 policy network 使用多层 convolution、zero padding、ReLU、softmax，但没有 pooling。

\begin{figure}[H]
    \centering
    \includegraphics[width=0.86\textwidth]{figures/fig16_alphago_no_pooling.jpg}
    \caption{AlphaGo 网络结构摘录：输入 $19\times19\times48$，第一层 $5\times5$ filters，后续 $3\times3$ filters；关键点是没有 pooling。画面依据：约 51:10--52:35。}
    \label{fig:alphago}
\end{figure}

\begin{importantbox}{架构的应用之道存乎一心}
    看到影像 CNN 常用 pooling，不代表所有 CNN 任务都该 pooling。每个架构元件都对应某种假设：locality、translation sharing、subsampling invariance。要把 CNN 用到语音、文字或围棋，必须重新检查这些假设是否符合任务。
\end{importantbox}

% =====================================================================
\section{CNN 的限制：不是天然处理缩放与旋转}
% =====================================================================

CNN 常被误解成天然具有各种视觉不变性。实际上，标准 CNN 主要内建的是局部性和位置共享，并不自动处理缩放、旋转。对人眼来说，同一只狗放大、缩小、旋转后仍是同一物件；但对网络输入向量来说，像素排列已经大幅改变。

\begin{figure}[H]
    \centering
    \includegraphics[width=0.86\textwidth]{figures/fig17_cnn_limitations.jpg}
    \caption{CNN 不天然 invariant to scaling and rotation；常需 data augmentation 或额外结构补足。画面依据：约 53:50--55:35。}
    \label{fig:limitations}
\end{figure}

这解释了为什么实际影像训练常做 data augmentation：随机裁切、缩放、旋转、翻转、颜色扰动等。Augmentation 的作用不是让资料变多这么简单，而是让模型在训练中看过同一物件的多种变形，从而学习到想要的不变性。

% =====================================================================
\section{本讲综合：CNN 是为影像写进网络里的先验}
% =====================================================================

这一讲的主线可以压缩成一句话：\textbf{CNN 把影像任务中合理的结构假设写进神经网络架构。}

\begin{itemize}
    \item 图片是三维 tensor，直接拉直接 fully connected layer 会产生大量参数。
    \item Receptive field 利用局部性：侦测小 pattern 不必看全图。
    \item Parameter sharing 利用平移共享：同一种 pattern 可在不同位置复用同一组参数。
    \item Filter 扫过图片的过程就是 convolution；多个 filters 产生多 channel feature map。
    \item Pooling 利用下采样不改变物件的直觉，但它不是所有任务都适合。
    \item CNN 的有效性来自任务先验与架构偏置匹配，而不是来自“卷积”这个名字本身。
\end{itemize}

\begin{knowledgebox}{和 P08/P10 的泛化观点连起来}
    P08 说模型复杂度越大，有限训练集越容易产生训练/真实落差。CNN 的做法不是盲目扩大模型，而是在影像问题上减少不必要的候选函数：不让每个 neuron 任意看全图，不让不同位置重复学习相同侦测器。只要这些限制符合任务，就能同时保留足够表达力，并降低过拟合风险。
\end{knowledgebox}

\end{document}
